% Generated draft: requires mathematical and source review.
\documentclass[11pt]{article}
\usepackage[margin=25mm]{geometry}
\usepackage{fontspec}
\setmainfont{Latin Modern Roman}
\usepackage{amsmath,amssymb,mathtools,unicode-math}
\setmathfont{Latin Modern Math}
\usepackage{enumitem,needspace,xurl,hyperref}
\setlength{\parindent}{0pt}
\setlength{\parskip}{5pt}
\setlength{\emergencystretch}{3em}
\newcommand{\N}{\mathbb{N}}
\newcommand{\Z}{\mathbb{Z}}
\newcommand{\Q}{\mathbb{Q}}
\newcommand{\R}{\mathbb{R}}
\newcommand{\C}{\mathbb{C}}
\newcommand{\F}{\mathbb{F}}
\newcommand{\A}{\mathbb{A}}
\newcommand{\PP}{\mathbb{P}}
\newcommand{\E}{\mathbb{E}}
\newcommand{\eps}{\varepsilon}
\newcommand{\dd}{\,\mathrm{d}}
\begin{document}
\section*{UnsolvedMath ID: 30003900}
\textbf{Source problem code:} catalogue entry
\subsection{Definitions and mathematical statement}
The supplied record describes a proposed relationship between two conjectures, but does not specify a complete implication. In particular, the meaning of its proposed \(\epsilon\)-strengthening is unresolved.

For clarity, the following definitions describe a conventional complex interpretation of the terminology. The source record does not specify the ground field; this interpretation therefore requires confirmation before it can become a definitive formulation. Let \(A\) be a complex abelian variety of dimension \(g\), and let \(\Theta\) be an effective divisor defining a principal polarization. The principally polarized abelian variety \((A,\Theta)\) is indecomposable if it is not isomorphic, as a polarized abelian variety, to a product of two principally polarized abelian varieties of positive dimension. Write \(\widehat A=\operatorname{Pic}^{0}(A)\). For a coherent sheaf \(\mathcal F\) on \(A\), set
\[
V^{i}(\mathcal F)=\{\alpha\in\widehat A:H^{i}(A,\mathcal F\otimes\alpha)\ne0\}.
\]
The generic-vanishing condition is
\[
\operatorname{codim}_{\widehat A}V^{i}(\mathcal F)\ge i\qquad\text{for every integer }i>0,
\]
with the empty locus assigned infinite codimension. For a closed subvariety \(X\subset A\), write \(\mathcal I_X(\Theta)=\mathcal I_X\otimes\mathcal O_A(\Theta)\), where \(\mathcal I_X\) is its ideal sheaf.

For a reduced irreducible complex subvariety \(X\) of dimension \(d\), the usual geometric nondegeneracy condition says that the kernel of restriction
\[
H^{0}(A,\Omega_A^{d})\longrightarrow H^{0}(X_{\mathrm{sm}},\Omega_{X_{\mathrm{sm}}}^{d})
\]
contains no nonzero decomposable form \(\omega_1\wedge\cdots\wedge\omega_d\), with \(\omega_j\in H^{0}(A,\Omega_A^{1})\). This definition does not resolve the record's missing assumptions on the scheme structure, irreducibility, or dimension range of \(X\).

For \(x\in\Theta\), if \(f\) is a local equation for \(\Theta\) in the regular local ring \(\mathcal O_{A,x}\), define
\[
\operatorname{mult}_{x}\Theta=\max\{m\ge0:f\in\mathfrak m_{A,x}^{m}\}.
\]
In the complex interpretation, a pure \(d\)-dimensional subvariety has minimal cohomology class when
\[
[X]=\frac{[\Theta]^{g-d}}{(g-d)!}\quad\text{in }H^{2(g-d)}(A,\Q).
\]
For a smooth projective complex curve \(C\), \(J(C)=\operatorname{Pic}^{0}(C)\), and \(W_d(C)\) denotes the image of the effective divisors of degree \(d\) in \(\operatorname{Pic}^{d}(C)\), transported to \(J(C)\) by a choice of translation. In the customary cubic-threefold interpretation, the Fano surface parametrizes lines on a smooth cubic threefold \(T\subset\PP^{4}_{\C}\), and its Abel--Jacobi image lies in the five-dimensional intermediate Jacobian \(J(T)\). The record does not explicitly supply this threefold hypothesis.

\textbf{Requested relationship, with unresolved qualifications.} Under the supplied hypotheses that \((A,\Theta)\) is indecomposable, \(X\subset A\) is geometrically nondegenerate, and \(\mathcal I_X(\Theta)\) satisfies generic vanishing, identify and quantify an additional strengthening sufficient to establish
\[
\forall x\in\Theta,\qquad \operatorname{mult}_{x}\Theta\le\frac{g+1}{2}.
\]
In the reverse direction, determine what consequences this multiplicity bound has for the proposed classification of \(X\) as a minimal-class subvariety, with the supplied target models being translates of \(W_d(C)\) in a Jacobian or of a Fano surface in a five-dimensional intermediate Jacobian, under suitable identifications of principally polarized abelian varieties.

The original wording calls these statements Conjecture 2 and Conjecture 4: an \(\epsilon\)-strengthening of Conjecture 2 is conjectured to imply Conjecture 4, and Conjecture 4 is expected to help prove important special cases of Conjecture 2. Their complete texts, the meaning and quantifiers of \(\epsilon\), and the intended special cases are absent. Thus neither the strengthened hypothesis nor a precise reverse implication can be stated from this record. Equality of cohomology classes and classification by the displayed geometric models are distinct assertions; their equivalence is not asserted here.
\subsection{Short English statement}
The record asks how a stronger generic-vanishing classification conjecture could imply the theta-divisor bound \(\operatorname{mult}_x\Theta\le(\dim A+1)/2\), and how that bound could assist classification in certain special cases. The strengthening, special cases, ground field, and several hypotheses are unspecified, so this remains a qualified statement of the requested relationship requiring human mathematical review.
\subsection{Sources}
\begin{itemize}
\item Supplied local UnsolvedMath registry record, permanent ID \(30003900\), exact source problem code \texttt{catalogue entry}, entitled ``Implications Among Conjectures on Principally Polarized Abelian Varieties,'' in the Oberwolfach Reports Open Problems collection. No individual proposer is supplied; the recorded proposed year is 2018.
\item Supplied report citation: \emph{Classical Algebraic Geometry (2019)}, Oberwolfach Reports, pp.~1983--2031, DOI \texttt{10.4171/owr/2018/33}. Supplied link: \url{https://doi.org/10.4171/owr/2018/33}. The citation's 2019 label, the DOI's 2018 component, and the recorded proposed year are retained as supplied.
\item Supplied literature source, labelled ``OWR report'': \url{https://ems.press/content/serial-article-files/46758?nt=1}.
\item Literature assessment dated exactly \texttt{2026-08-21}: the supplied curation labels the problem open and describes the exact record as an incomplete cross-reference between two conjectures rather than a complete standalone claim; it also reports that the referenced general classification remains unresolved. This is supplied dated curation, not an independently verified literature assessment. The background explicitly requests independent verification of time-sensitive status and cited claims.
\item The supplied statement status is \texttt{corrected\_verified}. Its verification note reports that the official report was checked and that the dangling conjecture numbers were replaced by generic-vanishing classification and theta-divisor multiplicity statements. This is a supplied provenance claim; no source visit or independent verification is claimed in this draft. The complete conjecture texts and a pinpoint location within the report were not supplied.
\end{itemize}
\end{document}
